\documentclass[11pt]{article}
\usepackage[margin=27mm]{geometry}
\usepackage{amsmath,amssymb,amsthm,booktabs,hyperref}
\newtheorem{theorem}{Theorem}
\newtheorem{lemma}[theorem]{Lemma}
\newtheorem{proposition}[theorem]{Proposition}
\newcommand{\Q}{\mathbb Q}\newcommand{\Z}{\mathbb Z}
\newcommand{\Tr}{\operatorname{Tr}}\newcommand{\Nm}{\operatorname{Nm}}
\newcommand{\rad}{\operatorname{rad}}
\title{Four scalar contacts and unsquared comparison scales}
\author{\(\Z\Q\) Diophantine Project: audit and continuation}
\date{19 November 2026}
\begin{document}\maketitle
\begin{abstract}
We audit the positive-factor and projective finite-algebra encodings supplied
in the two research bundles of Gabe Conant with Astra~6.2, dated 18 November 2026, and replace the scalar in
the general 25- and explicit 27-coordinate sources. Four finite contacts
suffice. A square-shifted input target separates integrality recognition
from scalar preparation. Unsquaring both nontrivial scales in the
27-coordinate source preserves its effective proof. The resulting degree
bounds are $87\,778\,394\,112$ and $335\,647\,080\,448$.
The latter is a completely specified member at $L=143$; the former still
has an ineffective sufficient fixed exponent. All defining statements
retain the precisely identified arithmetic premises of the source project.
\end{abstract}

\paragraph{Source credit.}
The supplied constructions and notes audited here are by Gabe Conant with
Astra~6.2, in the research bundles
\emph{H10-quartic-constructions} and
\emph{H10-eleven-variable-constructions}.
The audit, corrections, and further improvements stated in this paper are
subsequent project work. The original source files and their provenance
remain separately preserved.

\section{Scope, source systems, and result}

The base is the project repository cited below, commit \texttt{d3a8c22}.
The source manifest records its complete commit identifier and all file hashes.
We use the 25-coordinate source of v19, its unsquared-scale revision in
v28, its full auxiliary charts in v72, and its even-valuation tests in v75.
For 27 coordinates we use the three-scale v19 source and the full-family
modular estimate with small exceptional set recorded in v68. The notes of Gabe Conant with Astra~6.2
contribute the termwise-positive detector, the projective
two-pole encoder, and the improved effective height ledger.

The retained premises are the normalized group over $K=\Q(\sqrt2)$,
the relevant trace or effective abscissa-sum height theorem, the stated
modular contact-height theorem, the prepared native-pole converse and
its density assertion, and the compatible simultaneous prime-pattern
theorem for a fixed finite list of rational positions. This paper does
not reprove those foundations or make their audit implicit in a finite
coefficient check. It proves the transfers from them.

The four root roles $A,S,T,Z$ retain their square-abscissa charts and
have nonzero integral labels $n,a,t,b$. The five auxiliary roles use the
v72 full charts; their half-integral labels are integral at every consumed
odd prime. All old point, difference, logarithm, and comparison guards
are retained except the explicit automatic guards discussed below.
Write $c=|a-1|>0$, $B_0=\max(1,|n|,|a|,|t|,|b|)$, and
\[
 \Delta=t^2-a(a-2)tb+(a-2)^2b^2-1.
\]

\begin{theorem}\label{thm:main}
Under the corresponding retained premises there are integer-polynomial
presentations with the following bounds; every count includes the input.
\begin{center}\begin{tabular}{lrrr}
\toprule
Presentation&Coordinates&Degree at most&Support less than\\\midrule
General, sufficiently large odd $L$&25&$87\,778\,394\,112$&$2^{814}$\\
Fully specified, $L=143$&27&$335\,647\,080\,448$&$2^{929}$\\\bottomrule
\end{tabular}\end{center}
Their rational input projections are exactly $\Z$. For the 27-coordinate
member the coefficient one-norm is at most $2^{53113872449536}<2^{2^{46}}$.
The 25-coordinate receipt at $L=143$ is only an algebraic fixture, with
coefficient one-norm at most $2^{11446664560640}$; no assertion that this
particular exponent defines $\Z$ is made. The degree and support bounds
are uniform over sufficiently large odd fixed exponents.
\end{theorem}

\section{A scalar with four finite contacts}

Put $p_*=67$, and for odd $L\ge3$ define
\begin{align*}
 \beta&=1632X^2+2477,&\tau&=X^2+1,\\
 A(\beta)&=\beta(\beta-1),& B(\beta)&=(\beta-2401)(32\beta+343),\\
 T_L&=\bigl(1224(p_*^{(L+1)/2}+1)\bigr)^2,&q&=T_LA/B,\\
 \rho&=102^2p_*^L\frac{q}{(1+q)^2}
       =\frac{102^2p_*^LT_LAB}{(B+T_LA)^2}.
\end{align*}
Replace the first target test $L_1-\beta$ by $L_1-\tau$.
All sixteen comparison tests and six scalar components otherwise retain
their definitions and required formal orders. The scalar point remains
$S-P$ in the original normalized notation.

\begin{lemma}\label{lem:scalar}
For every rational $X$, $\rho$ is defined and positive,
$0<\rho<1/167$, and $v_{67}(\rho)=L$.
For integral $X$, $\rho$ has square reduced denominator, is a square in
$\Q_2,\Q_3,\Q_{17}$, and has respective valuations at least $10,6,4$.
Moreover
\[
 h(\tau)\le h(\beta),\qquad
 h(\rho)\le4h(\beta)+2L\log67+128.
\]
\end{lemma}
\begin{proof}
Since $\beta\ge2477$, both $A$ and $B$ are positive. Also
\[
 \frac BA<\frac{32\beta+343}{\beta-1}<33.
\]
Consequently $\rho<33/(144\cdot67)<1/167$.
The certificate enumerates all 67 values of $X$ modulo 67 and verifies
\[
 AB(B+1224^2A)\not\equiv0\pmod {67}.
\]
If $v_{67}(X)<0$, then $v_{67}(\beta)=2v_{67}(X)$ and
$A/B\equiv1/32$ after cancellation of their common valuation.
Both $1224^2$ and $32+1224^2$ are units modulo 67. Thus the claimed
valuation is exactly $L$ in both cases.

For any nonzero rational $q$, every negative valuation of
$q/(1+q)^2$ is even: if $v(q)\ne0$ its valuation is $|v(q)|$;
if $v(q)=0$ any negative valuation equals $-2v(1+q)$.
Multiplication by $102^2$ preserves parity, and the only remaining
prime, 67, has positive valuation. This proves the denominator assertion.

For integral $X$, $\beta\equiv13\pmod {32}$, so
$v_2(A)=v_2(B)=2$ and $AB/16\equiv3\pmod8$.
Since $67\equiv3\pmod8$ and $L$ is odd, the square class of
$67^LAB$ is trivial over $\Q_2$. The formula for $T_L$ gives
$v_2(T_L)\ge8$, hence $v_2(\rho)\ge10$.
Modulo 3 one has $AB\equiv1$ and $67\equiv1$, giving a square
class and $v_3(\rho)=6$. Modulo 17 one has $\beta\equiv12$,
$AB\equiv9$ and $67AB\equiv8$, a nonzero square; the valuation
of $\rho$ is at least four. All denominators in these square-class
calculations are squares.

If $X=u/v$ is reduced with $v>0$, the cancellation in $\beta$ is
\[
 g=\gcd(1632u^2+2477v^2,v^2)=\gcd(1632,v^2)\le1632.
\]
Its positive reduced numerator is therefore at least $u^2+v^2$.
As $\tau=(u^2+v^2)/v^2$ is reduced, $h(\tau)\le h(\beta)$.
Finally $\|A\|_1=2$, $\|B\|_1=900064$, and
$T_L\le 4\cdot1224^2\cdot67^{L+1}$. Put
$F_0=4\cdot1224^2\cdot67$. After factoring out $67^{2L}$,
the numerator and denominator coefficient one-norms of the degree-four
map in $\beta$ are bounded by
\[
 \max\{102^2F_0\cdot2\cdot900064,
              (900064+2F_0)^2\}<2^{128}.
\]
The standard elementary height estimate for a rational map proves the
last bound, with room to replace $128\log2$ by 128.
\end{proof}

The use of four finite contacts is essential to the improvement in the
circuit degree. The denominator contact of $\beta$ has even order
outside the fixed set containing the primes of 1632. At a finite odd
contact exactly one factor of $AB$ has odd valuation. Since $T_L$ is
a square, $v(q)$ is odd, and normalization makes the valuation of $\rho$
positive and odd. This argument includes primes dividing $T_L$; they
do not need to be added to the fixed exceptional set.

\section{Local tests and unsquared scales}

Use $\Theta=\{0,-1/2,-1,1/2399,-32/407\}$ and
\[
 \Theta_j=(4j+2+\Theta)\cup(4j+4+\Theta),\quad
 \delta_j=t_j+j+1,\quad
 R_j=\frac{\prod_{\theta\in\Theta_j}(t_j-\theta)}{\delta_j^{10}}.
\]
For two arguments set $t_0=\lambda$ and
$t_1=(1+\rho\mathcal T^2)\lambda$, where
$\mathcal T=\Tr\sum x(Q_i)$. This is the unsquared v28 scale.
For three arguments write $\sum x(Q_i)=U+\sqrt2S_1$, set $V=U+S_1$,
and use
\[
 t_0=\lambda,\qquad t_1=(1+\rho U^2)\lambda,\qquad
 t_2=(1+\rho V^2)\lambda.
\]
The outer squares in the old three-scale source are omitted.

For each detector keep the v75 reciprocal-abscissa numerators unsquared
and all other raw tests and weights unchanged. Replace every factor
$1+a$ by $1+a^2$, including the normalized scalar factors.
At a consumed odd valuation each nonzero $a$ has odd valuation, so
\[
 v(1+a^2)=2\min(0,v(a))\le0.
\]
The negative-odd pole obstruction therefore implies at least the old
local depths. At a freshly selected prime the successful factors are
one modulo the prime, and the opposite inverse factors are zero or one.
The trace remains a unit. At fixed finite and real preparation places
the old small raw tests give the same limiting trace two. Zero raw
tests contribute one throughout.

For clarity, the unsquared scales preserve the scalar-unit argument.
At a prime with odd $v_p(\rho)$, each $\rho U^2$ has odd valuation;
therefore a nontrivial scale is either congruent to one or has negative
valuation. If $\lambda$ is nonintegral or has positive valuation,
$R_0$ is a unit at the relevant nonexceptional primes. Otherwise
either some $t_j$ is nonintegral, making $R_j$ a unit, or all $t_j$
have the same residue as $\lambda$. The old exclusion of common
contact/pole residues applies. The two sets are disjoint modulo 67,
as is independently checked. Thus the fixed-prime extraction works
even though 67 lies below the old large technical cutoff. The good
formal-logarithm estimate of v68 supplies the same relative order four
at this prime. No new modular theorem is needed for the two-scale
ineffective argument, whose original large-cutoff height estimate may
be retained.

Consequently, with the relevant excluded finite sets, the old scalar
and mixed-contact lemmas give
\[
 67^{2L+1}M_\beta^3\mid c,\qquad
 M_j^3\mid\Delta c,\qquad M_j\mid Nc,
 \quad N=\operatorname{num}(\tau-n),
\]
whenever a residual is nonzero. The first root label is integral;
there is no half-integral-label shortcut here.

\section{Effective soundness for three arguments}

We reproduce the numerical ledger so that unsquaring the scales is
paid for explicitly. Put
\[
 S=\{2,3,5,7,11,13,17,37,41,2399\},\quad
 \Gamma=S\cup\{19,23,457,4799\}.
\]
The differences among $0,1,2401,-343/32$ add no prime to $\Gamma$.
The full-family modular premise is used in exactly the form
\[
 h(q)<C_*R^2E(R)\log R,\quad
 C_*=\frac{74712001}{10^6}2^{68}3^{16}5^{16}7^4,
 \quad E(R)=\prod_{p\mid R}(1+1/p).
\]
The independent integer ledger verifies the scalar bound
$h(\beta)<2^{342}M_\beta^3$ and the three contact bounds
\[
 h(t_j)\le K M_j(199+\tfrac12\log M_j)^2,
 \qquad K=2^{551}.
\]
For the latter, at a prime in the enlarged $j$th exceptional set,
both five-contact radicals can contain that prime only when contacts
in opposite blocks collide. Counting that prime twice precisely in
this case gives $K_j<2^{398}$ and $C_*K_jE_j<2^{550}$.
Taking the geometric mean of the two modular estimates gives the
displayed bound. The checker reconstructs every finite prime set and
every multiplicity, not just the resulting powers of two.

The effective abscissa-sum premise gives
$B_0^2\le2h(U)+2h(V)+82$. The new scales satisfy
\[
 h(1+\rho U^2)\ge2h(U)-h(\rho)-\log2
\]
and its analogue for $V$. Hence
\[
 B_0^2\le2h(t_0)+h(t_1)+h(t_2)+2h(\rho)+84.
\]
Keeping the weaker scalar coefficient bound $15h(\beta)$, a nonzero
residual $N$ of order $k=3$ or $1$ now gives
\[
 B_0^2\le4K|Nc|^{1/k}
       (199+\tfrac1{2k}\log|Nc|)^2
       +30h(\beta)+4L\log67+340.                 \tag{*}
\]
The factor four, rather than two, is the exact cost of unsquaring.

For every odd $L\ge143$, the certificate verifies
\[
 67^{2L+1}\ge120\cdot2^{342},\qquad
 c-1\ge2^{1727}.
\]
Thus $h(\beta)<c/120$ and $4L\log67\le2\log c$.
The last three terms of (*) are less than $B_0^2/2$.
If $\Delta\ne0$, the source integer estimate
$|\Delta c|\le10B_0^5$ would imply
\[
 B_0^{1/3}<18K(200+\tfrac56\log B_0)^2.
\]
The quotient of the two variable factors is increasing for $B_0\ge1$.
At $B_0=2^{1727}$ the reverse inequality follows from the exact test
\[
 2^{1727-3\cdot551}12^6>18^3(2400+7\cdot1727)^6
\]
and $\log2<7/10$. Therefore $\Delta=0$, and the retained elementary
Lucas descent gives $\log B_0>c-4$.

If $N=\operatorname{num}(\tau-n)\ne0$, then
$|N|\le2e^{h(\tau)}B_0\le2e^{h(\beta)}B_0$.
Writing $s=\log B_0$ and $t=h(\beta)+\log c$, the $k=1$ case of
(*) yields
\[
 s\le\log(16K)+t+2\log(200+(s+t)/2).
\]
Using $\log(1+x)\le x$ and the exact inequality
$200(551+4+16)<199\cdot574$ gives
\[
 s<2t+574\le(31/60)c+574<c-4,
\]
a contradiction. Thus $\tau=n\in\Z$. Since
$\tau=(u^2+v^2)/v^2$ is reduced, this implies $v=1$ and $X\in\Z$.

For two arguments the same transfer retains the v28 inequality with
unspecified fixed constants, scalar loss $c^{7/9}$, and comparison
exponents $7/18$ and $7/6$. The stronger scalar height bound and
$h(\tau)\le h(\beta)$ preserve both contradictions. The fixed-prime
lower bound tends to infinity with $L$. Its sufficient threshold remains
ineffective because the retained trace-height constant is ineffective.

\section{Completeness and guards}

Fix an integer $X$ and put $n=X^2+1\ge1$. Choose the scalar reference
neighborhood, then the Lucas labels, then all root points and scales,
in the original order. No proof step concerning the first root requires
that $n$ be odd: it must be nonzero, and its finitely many prime divisors
are avoided before selecting comparison primes. The other Lucas parity
conditions are unchanged. The auxiliary points may be chosen in the
original $2E(K)$ subfamily.

Lemma~\ref{lem:scalar} gives every scalar square-class, denominator,
depth and real-size condition used in preparation. The unsquared
scales are positive rational numbers. In the prime factorization of
each balanced ten-contact function the scale occurs to exponent ten;
its contribution is a square whether or not the scale itself is a
square. All finite cutoffs and denominator clearing are fixed before
prime selection. If scaled positions from different blocks coincide,
the prime list records that position once and retains every incidence,
as in v19. Within each block the contacts are simple; across blocks the
same successful elliptic comparisons are prepared. Thus repeated
positions impose compatible requirements. The simultaneous prime-pattern theorem and compatible
index refinement therefore supply the old integer witnesses. Positive
factors have the same residues and limiting traces, so the actual pole
inputs satisfy the same restricted converse. Its dense choices avoid
the finitely many remaining pole conditions. This includes $X=0$:
its first root label is one, and no separate input factor is needed.

The finite-algebra guard can omit the atoms in $\rho$ and the scale
numerators which are automatically nonzero. Indeed all four scalar
factors are positive for every rational input, $B+T_LA>0$, and every
scale is $1+\rho r^2>0$ for a rational base-coordinate expression $r$.
The checker proves that these expressions involve no finite generator.
Their positivity therefore precedes, and does not depend on, rational
character recovery. The later positive-detector guard deletions are
the audited two-stage argument of the notes of Gabe Conant with Astra~6.2. They are used
only after the remaining protected denominators and a rational character
have been obtained.

\section{Finite-algebra specifications and independent checks}

All point algebras, their integral bases and quadratic envelopes are
the pinned v28/v30/v72/v75 ones. Zero terminal pole ordinates are
allowed. For three poles $q_0,q_1,q_2$ over a point character use
$E=q_0^3+2q_1^3+4q_2^3$ as the terminal sign separator. In a common
Galois 2-extension, at a prime over 2 the nonzero summands of any
signed difference have pairwise different valuations modulo three.
They cannot cancel. The bounded point-prefix construction then
separates distinct point characters and gives a nonzero integral
separator $\Omega$. The polynomial is $\det(I-YG\,\mathrm{Mult}(\Omega))$.

For two poles use the audited projective terminal value
\[
 \frac18\left(x^3+\frac{z}{1+z}\right),\qquad
 x=\frac{q_1}{2M_1},\quad z=\frac{q_2}{2M_2},\quad |x|,|z|\le\tfrac12.
\]
Once the point character is rational, $x^2=d,z^2=e$ are rational.
If the displayed value is rational and a root is irrational, both
roots must lie in the same quadratic field. With $z=rx$, cancellation
of their irrational parts yields $d^2r^2-r-d=0$ and makes
$1+4d^3$ a square. The curve $y^2+y=d^3$ has only $d=0$ among its
affine rational points, contrary to irrationality of $x$.
This classical exponent-three fact is proved by the Fermat reduction
in \texttt{PROJECTIVE.md} by Gabe Conant with Astra~6.2; it is also the rank-zero,
order-three curve 27.a4 in LMFDB. Appending the bounded point prefix
therefore gives a numerator/denominator pair $N_1,D_1$ for which
\[
 \det(\mathrm{Mult}(D_1)-YG\,\mathrm{Mult}(N_1))
\]
has a rational zero exactly when the source has a rational character.
The denominator is nonzero on every geometric character. At $G=0$
the determinant is its nonzero norm. Nilpotent specializations cause
no problem: multiplication on each Artin local factor has the residue
character value on its diagonal. One may divide by the positive
integer content to obtain a primitive polynomial without increasing
any stated bound.

There is one further guard reduction valid for both sources. In a root
chart write $x=g^2$ and $j^2=g^4-D$, with $D=11-6\sqrt2$.
The norm-root square $\Nm(g^4-D)$ is nonzero at every rational base:
otherwise $g^4=D$, and $(\Nm g)^4=49$, impossible over $\Q$.
Its norm guard is therefore automatic before extracting a character.
For a full auxiliary chart, the retained ordinate-norm guard makes
$\Nm(x^3-Dx)$ nonzero, which forces $\Nm(x)\ne0$.
These facts restore point etaleness before using the integral point
separator. After a rational character has been extracted, $g=0$ would
give $j^2=-D$, impossible because $D$ is totally positive.
Thus the root abscissa norm is recovered at this later stage.
The integral multiplication algebra exists on all the removed loci;
raw rational denominators are evaluated only after their recovery.
For four root and five auxiliary charts this removes thirteen norm
atoms of total guard degree 50. It does not assert positivity at
nonrational characters and is independent of the detector deletions.

The rank is $2^{20}$ for 25 coordinates and $2^{21}$ for 27.
The retained guard degrees are respectively $8375$ and $14698$.
For 25 the projective numerator and denominator operator degrees are
both $75336$; hence the bound is
$2^{20}(1+8375+75336)=87778394112$.
For 27 the separator degree is $145350$, giving
$2^{21}(1+14698+145350)=335647080448$.
The support counts use the exact binomial difference for the total
degree and the encoder degree at most the algebra rank.

\texttt{check\_scalar\_arithmetic.py} is a standard-library checker
for all residues, exceptional prime sets, scalar constants and threshold
inequalities. It imports neither the producer nor a CAS.
\texttt{check\_scalar\_circuits.py} independently reconstructs the
scalar, target, all scales, all contact products and weights; reconciles
every unchanged chart and test with the pinned source; and replays all
44 or 66 positive detector factors, 411 or 534 denominator records,
square extractions, envelopes, guard degrees, and coefficient majorants.
The finite checks certify those algebraic assertions; the preceding
proof supplies their arithmetic meaning under the retained premises.

\begin{thebibliography}{9}\raggedright
\bibitem{sources} \(\Z\Q\) Diophantine Project, pinned repository above:
25- and 27-coordinate v19 papers; unit-norm v28; subalgebra-norm v30;
terminal-root v40; compact quartic v68; full auxiliary v72;
even-valuation v75. All sources accompany the reproducibility bundle.
\bibitem{coauthor} Gabe Conant with Astra~6.2, research bundles, 18 November 2026:
\emph{H10-quartic-constructions} and \emph{H10-eleven-variable-constructions};
in particular \texttt{POSITIVE.md}, \texttt{PROJECTIVE.md}, and the
effective \texttt{PROOF.md} and \texttt{COMPACT.md}.
\bibitem{lmfdb} LMFDB, elliptic curve 27.a4,
\url{https://www.lmfdb.org/EllipticCurve/Q/27/a/4}, checked 19 November 2026.
\end{thebibliography}
\end{document}
